\documentclass{article}
\usepackage[margin=25mm]{geometry}
\usepackage{amsmath,amssymb}
\usepackage{xurl}
\usepackage[hidelinks]{hyperref}
% Shared commands for notes in docs/paper-gaps/.

\DeclareUnicodeCharacter{2080}{\ifmmode{}_{0}\else\textsubscript{0}\fi}
\DeclareUnicodeCharacter{2081}{\ifmmode{}_{1}\else\textsubscript{1}\fi}
\DeclareUnicodeCharacter{2082}{\ifmmode{}_{2}\else\textsubscript{2}\fi}
\DeclareUnicodeCharacter{2099}{\ifmmode{}_{n}\else\textsubscript{n}\fi}

\newcommand{\Lean}{\textsc{Lean}}
\ExplSyntaxOn
\NewDocumentCommand{\leanid}{m}{
  \ifmmode
    \textsf{\detokenize{#1}}
  \else
    \group_begin:
    \urlstyle{sf}
    \tl_set:Nn \l_tmpa_tl {#1}
    \tl_replace_all:Nnn \l_tmpa_tl {\_} {_}
    \exp_args:NV \path \l_tmpa_tl
    \group_end:
  \fi
}
\ExplSyntaxOff
\providecommand{\tr}{\operatorname{tr}}
\newcommand{\tnleanIssueBase}{https://github.com/LionSR/TNLean/issues/}
\newcommand{\tnleanPrBase}{https://github.com/LionSR/TNLean/pull/}
\newcommand{\tnleanDocBase}{https://sirui-lu.com/TNLean/docs/}
\newcommand{\ghissue}[1]{\href{\tnleanIssueBase#1}{GitHub issue \##1}}
\newcommand{\ghpr}[1]{\href{\tnleanPrBase#1}{PR \##1}}
\newcommand{\doclink}[1]{\href{\tnleanDocBase#1.html}{\path{#1}}}

% Dirac notation and the identity operator, as in blueprint/src/macros/common.tex.
% \providecommand keeps note-local definitions authoritative.
\usepackage{dsfont}
\providecommand{\ket}[1]{|#1\rangle}
\providecommand{\bra}[1]{\langle #1|}
\providecommand{\braket}[2]{\langle #1 | #2 \rangle}
\providecommand{\ketbra}[2]{|#1\rangle\!\langle #2|}
\providecommand{\Id}{\mathds{1}}
\providecommand{\one}{\mathds{1}}
\providecommand{\C}{\mathbb{C}}
\providecommand{\MN}[1]{M_{#1}(\C)}
\providecommand{\MD}{M_D(\C)}

% Verdict marker: \gapnote{<kind>}{<status>}. Kind is the policy
% classification (clarification, local-correction, scope-restriction,
% unfaithful, false-source, open-gap); status is open, wip (elimination
% underway), resolved, or historical. Machine-read by the site index;
% prints a verdict line.
\newcommand{\gapnote}[2]{\par\noindent\textbf{Verdict:} #1 (#2).\par}

\providecommand{\gapnote}[2]{\par\noindent\textbf{Verdict:} #1 (#2).\par}
\title{Ground-Energy Dependence in the Sampled Projection of Landau, Vazirani and Vidick}
\author{}
\date{October 29, 2026}
\begin{document}
\maketitle
\gapnote{false-source}{open}

\section{The preprint and the published correction}

The sampling construction in arXiv:1307.5143v1~\cite{Landau2015Polynomial},
Section~2.3.1, Corollary~2.14 (\path{c:agsp}; \path{algo53.tex},
lines~628--649), uses
\begin{align}
    A   & =\bigl(C(I-H/n)\bigr)^m,                   \\
    C   & =(1-\varepsilon_0/n)^{-1},
    \label{eq:lvv_sampling_power}                    \\
    K   & =\frac{C^m}{\ell}\sum_{j=1}^{\ell}P_{I_j}, \\
    P_I & =\prod_{k=1}^{m}P_{i_k}.
    \label{eq:lvv_sampling_average}
\end{align}
Here $P_i=I-H_i$. There are $n-1$ physical interactions; the missing label is
read as $P_n=I$. The word $I_j$ is uniform on $\{1,\ldots,n\}^m$. The proof asserts
$C^{2m}=n^{O(1/\varepsilon)}$ for the power
$m=\Theta((n/\varepsilon)\log(q(n)/c_\varepsilon))$ required by
equation~\path{e:cutdown}. Its hypotheses impose no bound on
$\varepsilon_0$ independent of $n$.

The published article~\cite{Landau2015Published}, Methods, states the
complexity exponent as
\begin{align}
    c(d,\varepsilon_0,\varepsilon)
    =2^{O(\log^3d/\varepsilon)}+O(\varepsilon_0/\varepsilon).
    \label{eq:lvv_sampling_published_exponent}
\end{align}
The Methods statement of Proposition~1, author PDF page~6, retains the
ground energy in the sampling complexity. The normalization estimate below
makes this dependence explicit. The abstract advertises the uniform
polynomial case with constant ground energy.
The published theorem treats $\varepsilon_0$ as a complexity parameter;
it does not impose an unprinted bounded-energy hypothesis. Thus the omission
of $\varepsilon_0$ in the preprint was corrected in publication. The example below concerns that omitted dependence, and does
not contradict the published theorem.

\section{A commuting chain with zero local minima}

Let $n=2r+1$ with $r\ge1$, and let $N_j=|1\rangle\langle1|_j$ on a chain
of $n$ qubits. For $1\le k\le r$, set
\begin{align}
    H_{2k-1} & =\frac34N_{2k}+\frac18N_{2k-1},     \\
    H_{2k}   & =\frac12(I-N_{2k})+\frac18N_{2k+1}.
    \label{eq:lvv_sampling_chain}
\end{align}
Each term is supported on its prescribed adjacent pair, has smallest
eigenvalue zero, and satisfies $0\le H_i\le(7/8)I$. All terms commute.
Their sum is
\begin{align}
    H=\frac r2I+\frac14\sum_{k=1}^{r}N_{2k}
    +\frac18(N_1+N_n)+\frac14\sum_{k=1}^{r-1}N_{2k+1}.
    \label{eq:lvv_sampling_hamiltonian}
\end{align}
Consequently $\Gamma=|0\rangle^{\otimes n}$ is the unique ground vector,
$\varepsilon_0=r/2$, and the gap is exactly $1/8$, attained by flipping
either endpoint. This is a nondegenerate example even if zero local
minimum is imposed as an additional convention.

\section{Failure of the claimed normalized operator-norm estimate}

Put $p=r/n$, so $1/3\le p<1/2$. On $\Gamma$, the operators $P_i$ have
eigenvalue $1$ for odd $i$ and for the dummy label, and eigenvalue $1/2$
for even $i$. For a uniform word of length $m$, let $B$ count the even
labels. Then $B$ has the binomial distribution with parameters $m,p$, and
the eigenvalue of $C^mP_I$ on $\Gamma$ is $Z=C^m2^{-B}$, with
$C=(1-p/2)^{-1}$. In particular,
\begin{align}
    \mathbb EZ=\bigl(C(1-p/2)\bigr)^m=1.
    \label{eq:lvv_sampling_word_law}
\end{align}
Its square-root moment is
\begin{align}
    \mathbb E\sqrt Z & =\alpha_p^m,                        \\
    \alpha_p         & =\frac{1-p+p/\sqrt2}{\sqrt{1-p/2}}.
    \label{eq:lvv_sampling_fractional_moment}
\end{align}
Writing $t=1-1/\sqrt2$, direct expansion gives
\begin{align}
    \alpha_p^2
    =1-\frac{p(1-p)t^2}{1-p/2}
    \le 1-\frac1{96}=\frac{95}{96}.
    \label{eq:lvv_sampling_moment_decay}
\end{align}
Indeed $p(1-p)\ge1/6$, $t\ge1/4$, and $0<1-p/2\le1$.

For $\ell\ge1$ sampled words, let
$Y=\ell^{-1}\sum_{j=1}^{\ell}Z_j$, so $K\Gamma=Y\Gamma$.
The inequality $\sqrt{\sum_jZ_j}\le\sum_j\sqrt{Z_j}$ and Markov's
inequality imply
\begin{align}
    \Pr(Y\ge1/2)
     & \le \sqrt2\,\mathbb E\sqrt Y
    \le\sqrt{2\ell}\,\alpha_p^m
    \le\sqrt{2\ell}\left(\frac{95}{96}\right)^{m/2}.
    \label{eq:lvv_sampling_failure}
\end{align}
Since $A\Gamma=\Gamma$, the event $\|K-A\|\le1/n$ implies
$Y\ge1-1/n\ge1/2$. For every fixed polynomial bound on $\ell$ and every
$m\ge cn$ with $c>0$ fixed, the right side of
(\ref{eq:lvv_sampling_failure}) tends to zero. In particular, the
prescribed power of order $n\log n$ cannot have the claimed success
probability with a polynomial number of uniform samples.
This concerns the stated normalization and additive operator-norm estimate.
It does not say that every normalized output vector of this commuting
example is poorly approximated: a nonzero scalar multiple of the ground
vector has the same normalized direction.

The bare upper bound $m=O(n\log n)$ also permits small powers, including
$m=0$. The contradiction is to the sampled approximate projection at
the powers required by equation~\path{e:cutdown}, not to an
upper-bound statement interpreted without that requirement.

\section{Quantitative repair and formalization boundary}

The normalization already displays the correct energy dependence. Put
$a=n-\varepsilon_0\ge1$, choose $Q\ge1$, and take
$m=\lceil a\log Q/\varepsilon\rceil$. Then
\begin{align}
    \log C
     & =\log(1+\varepsilon_0/a)\le\varepsilon_0/a, \\
    C^{2m}
     & \le C^2Q^{2\varepsilon_0/\varepsilon}
    \le n^2Q^{2\varepsilon_0/\varepsilon}.
    \label{eq:lvv_sampling_energy_dependence}
\end{align}
Here $\varepsilon_0\ge0$, and the ceiling costs the factor $C^2$.
The usual matrix concentration estimate can therefore give a polynomial
sample count on families with bounded $\varepsilon_0$, or an
energy-dependent exponent as in the published statement. A stronger
concentration theorem cannot give the stated normalized operator-norm
approximation with that energy-independent polynomial count on
(\ref{eq:lvv_sampling_chain}).

The uniform polynomial construction for unrestricted gapped chains
requires a different approximate projection or a derived energy
reduction. A different finite Gaussian construction for arbitrary ground
energy is given independently by Huang~\cite[Section~5.4]{Huang2015GroundState},
based on locality estimates. The energy-dependent construction must retain
its ground-energy exponent.
The Gaussian construction described below supplies the successive
polynomial existence theorem without a bound on the ground energy.
It does not establish the preprint's estimate for its uniform sampler.

\section{The actual high-probability event}

The probability estimate itself is preserved in
\path{TNLean/MPS/AreaLaw/ViableSet/SampledProjector.lean}.
The underlying law is the actual product of uniform letters on
$\{0,\ldots,n-1\}^{\ell\times m}$, with the last label denoting the
identity padding letter. For fixed physical ground data, $m\ge0$ and
$\eta>0$, put $R=C^m+1$ and choose
\begin{align*}
 \ell&=\left\lceil
   \frac{2R^2+2R\eta/3}{\eta^2}\log(8q^n n^3)
   \right\rceil+1,\\
 \kappa&=\left\lceil 2m/n+2\log(4\ell n^4)\right\rceil.
\end{align*}
The joint event $\|K-A\|<\eta$ and at most $\kappa$ occurrences of every
letter in every sampled word has probability at least $1-1/(2n^3)$,
hence at least the source's $1-1/n^3$. The matrix-accuracy failure and
all-word multiplicity failure are each bounded by $1/(4n^3)$ under this
same product law.

This assertion specifies the distribution and the joint event. It is
stronger than existence of one successful family, and is not inferred
from the separate existence-only or bracketed sampling constructions.
Its explicit count retains the ground-energy dependence in
(\ref{eq:lvv_sampling_energy_dependence}), consistently with the published
Methods correction. It does not recover the preprint's energy-independent
exponent or choose unknown-energy trials. The distinction between this
probability theorem and a trial-energy algorithm remains essential.

\section{Overall scaling and an unknown energy}

The scalar $C^m$ is not itself an obstacle to using the left-factor
construction without an energy input. The published Methods discussion before Proposition~1,
author PDF page~6, explicitly observes that an overall scalar is immaterial after
normalizing the output vectors. The same point occurs in
\path{algo53.tex}, line~614. One may use the unscaled sum
$\ell^{-1}\sum_jP_{I_j}$; multiplication of all factors by one nonzero
scalar preserves their spans and Schmidt ranks.

What remains energy-dependent in the present formal statement is the
selection of its word length, relative error and sample budget.
A known upper bound on the energy, or a proved sequence of trial bounds
with a valid acceptance and termination argument, would address a different
question from computing the immaterial scalar. Such a bound cannot be
inferred merely by subtracting the global ground energy while retaining
positive original local contractions.

Both source versions use the power of the full original nearest-neighbour
Hamiltonian. Its left and right block sums occur in the size-trimming
optimization; they do not replace the original letters in this sampler.
Subtracting their two separate minima gives a residual ground energy in
$[0,1]$, but the shifted block operators are nonlocal and can have norm of
order the chain length. Their internal-cut ranks are not controlled by the
original-letter estimate. Substituting them would require new norm and
spatial-rank arguments.

\section{Uniform successive families and arbitrary accuracy}

Fix the local dimension $q\ge2$ and a positive lower bound $\Delta$ for
the spectral gap. The Gaussian energy-mesh construction supplies a
sequence of viable families whose size and individual bond dimensions
are bounded by one polynomial in the chain length. The coefficients and
degree are chosen before the length and interaction, and the families
are chosen before any ground energy or ground vector. The theorem is
\leanid{TNLean.MPS.AreaLaw.exists_uniform_positive_chain_successive_approximation}.
Its proof performs extension, size trimming, independent bond trimming
and Gaussian error reduction at each cut. The final approximation lies
in the span of the family produced at the last cut.

Arbitrary accuracy follows by adjoining $2R^2$ sites with product-state
penalties, for a prescribed integer $R\ge1$. The extended chain has
length $N=L+2R^2$ and gap at least $g=\min\{\Delta,1\}$. A normalized
column of the stage-$L$ witness, with its phase chosen, gives an
original-chain matrix product state $\psi$. There are natural numbers
$P,d$, depending only on $q,\Delta$, such that
\begin{align}
    \|\psi-\Gamma\|^2 &\le R^{-2}, &
    D(\psi) &\le P(L+2R^2+1)^d.
    \label{eq:lvv_gaussian_arbitrary_accuracy}
\end{align}
The same polynomial bounds the cardinalities and individual bond
dimensions at every stage of the extended construction. This is
\leanid{TNLean.MPS.AreaLaw.exists_uniform_positive_chain_successive_accuracy}.
The extension by product states is a consequence of the construction,
not a step claimed to occur in the cited papers.

These results use the corrected trimming radius and the constant-cut
bond described in \path{lvv13_constant_bond.tex}. They require no bound
on the ground energy. They do not repair the false normalized estimate
for the original uniform sampler, establish the sharper printed
dependence on the inverse gap, or assert a computational running time.
The source-level construction and complexity comparisons remain
tracked by issues \#9822 and \#9823 under \#9832.

\bibliographystyle{alpha}
\bibliography{references}
\end{document}
